\section{Historical return evidence}\label{sec:baseline_return_predictability}
\subsection{Portfolio sorts}
We reproduce the Annual Review program with the original scores, sign restrictions, quantile cutoffs, position limits and eligibility rules. The long leg contains up to 50 positive-signal stocks in the upper 30 percent of the daily distribution; the short leg contains up to 50 negative-signal stocks in the lower 30 percent. Both legs must exist. EW portfolios give equal weight within each leg; VW portfolios use the original market-capitalization weights.

Table~\ref{tab:horse_race} reports gross ESI returns of 5.31 basis points per day with equal weights ($t=5.95$) and 5.26 with value weights ($t=3.57$). Arithmetic annualized returns are 12.85\% and 12.72\%. Investor-tone portfolios earn 4.43 and 4.38 basis points; manager-tone portfolios earn 2.48 and 2.75. ESI has the largest point estimate among these three signals under both weights. Separate significance levels and unequal candidate pools do not establish that differences are statistically significant.

The domestic CH4 model includes market, size, value and sentiment factors \citep{liu2019}. The historical q5 comparison uses the U.S. Global-q augmented q-factor series \citep{hou2021}; it is not a Chinese CSMAR factor set. ESI alphas are positive under both benchmarks. The original signal-date factors and daily intervals differ from next-record intraday returns, so these are diagnostic historical alphas. The availability comparison separately reports domestic CH4 and domestic FF3 over aligned holding dates.

The original program excludes stocks according to their next recorded closing-limit status, which is unavailable at signal-date portfolio formation. Together with full-period scales and retrospective reply assignment, this prevents interpretation as attainable trading profits. The short leg is theoretical, and returns exclude trading, borrowing and market-impact costs. Recomputed values replace old printed figures without changing parameters to make them coincide.
\begin{table}[H]
\centering\footnotesize
\caption{Historical O2C Portfolio Performance}\label{tab:horse_race}
\begin{threeparttable}
\renewcommand{\arraystretch}{1.12}
\begin{tabular*}{\textwidth}{@{\extracolsep{\fill}}lrrrrr@{}}
\toprule
Signal & Mean (bps) & Annual (\%) & Sharpe & CH4 $\alpha$ & U.S. q5 $\alpha$ \\
\midrule
ESI (EW) & 5.31 (5.95) & 12.85 & 1.54 & 5.17 (5.74) & 5.42 (6.06) \\
Investor tone (EW) & 4.43 (4.22) & 10.73 & 1.11 & 4.73 (4.55) & 4.60 (4.34) \\
Manager tone (EW) & 2.48 (2.83) & 6.00 & 0.71 & 2.36 (2.65) & 2.46 (2.80) \\
ESI (VW) & 5.26 (3.57) & 12.72 & 0.91 & 5.16 (3.46) & 5.40 (3.67) \\
Investor tone (VW) & 4.38 (2.70) & 10.59 & 0.69 & 4.73 (2.88) & 4.51 (2.79) \\
Manager tone (VW) & 2.75 (1.79) & 6.66 & 0.45 & 2.73 (1.77) & 2.69 (1.75) \\
\bottomrule
\end{tabular*}
\begin{tablenoteblock}
Gross long-minus-short returns; 3,642 dates for each row. Parentheses contain HAC(5) $t$ statistics. Up to 50 positive-signal stocks in the upper 30 percent form the long leg; up to 50 negative-signal stocks in the lower 30 percent form the short leg. Both legs must exist. Annualization is the historical arithmetic convention, $242\bar r\times100$, with Sharpe scaled by $\sqrt{242}$. Candidate pools are signal-specific, not a common-stock-pool horse race. Next-record O2C excess returns, future closing-limit exclusions, full-period text scales, and signal-date factors reproduce the original implementation. CH4 is domestic; q5 is the U.S. Global-q series. The intraday outcome and daily factor intervals differ. These are historical diagnostic alphas, not implementable profits. All figures are recomputed values rather than copied PDF numbers.
\end{tablenoteblock}
\end{threeparttable}
\end{table}


\subsection{Return regressions}
The historical baseline estimates
\begin{equation}
y^{O2C}_{i,t+1}=\alpha_i+\delta_t+\beta Z^{ESI}_{i,t}+\boldsymbol c_{i,t}'\boldsymbol\eta+\varepsilon_{i,t+1},\label{eq:rpr}
\end{equation}
where $Z^{ESI}$ uses frozen full-regression-pool moments and $\boldsymbol c$ contains the eleven controls in Table~\ref{tab:variables}. Firm and question-date effects absorb additive differences across firms and dates, and errors are clustered by both dimensions. The index $t+1$ follows the original next-record convention; it need not be the first exchange date after all information becomes public.

Table~\ref{tab:panel_horse_race} reports 1.78 basis points per ESI standard deviation for the q5 residual ($t=5.48$, $N=920{,}584$). China CH4 and raw-excess-return estimates are also positive, at 1.78 and 1.79 basis points. Different outcome-availability samples are stated explicitly. These are conditional one-day associations, not joint tests against every ESI component. The old multi-horizon O2C and Fama--MacBeth tables remain archived because their printed values have not been reconciled with a complete re-estimation.
\begin{table}[H]
\centering\footnotesize
\caption{Historical One-Day O2C Regressions}\label{tab:panel_horse_race}
\begin{threeparttable}
\renewcommand{\arraystretch}{1.12}
\begin{tabular*}{\textwidth}{@{\extracolsep{\fill}}lrrrr@{}}
\toprule
Dependent variable & ESI (bps) & $t$ & Firm days & Firms \\
\midrule
U.S. q5 residual & 1.78 & 5.48 & 920,584 & 3,593 \\
China CH4 residual & 1.78 & 5.60 & 953,080 & 3,594 \\
Raw O2C excess return & 1.79 & 5.60 & 953,081 & 3,595 \\
\bottomrule
\end{tabular*}
\begin{tablenoteblock}
Each row is a separate regression with the eleven baseline controls, firm effects, and question-date effects; standard errors are clustered by firm and date. ESI is standardized once using the original full regression pool. Residual outcomes use full-sample firm loadings and original signal-date factors. Controls and outcomes determine each row\textquotesingle s sample; these are not common-sample factor comparisons. Positive one-day associations do not establish independent text information or real-time forecasting. The unverified multi-horizon O2C and Fama--MacBeth tables remain in the historical archive.
\end{tablenoteblock}
\end{threeparttable}
\end{table}


\subsection{Response lag}
We classify each original company-question-date cohort by the maximum natural-day lag among recorded replies. Same-day means all valid replies occur on the question date. Other groups are 1 day, 2 days, 3--5 days, and more than 5 days. There is no restriction that all replies occur on one common reply date. For example, replies to a Day~1 cohort arriving on Days~2 and~4 imply a maximum lag of 3 days; its retrospective ESI remains on Day~1.

Among the 865,309 firm days in the common grouped regression sample, the same-day share is 11.56\%. Portfolio and regression populations differ because regressions require additional controls; separate distributions and exclusions are preserved. This sample differs from the later public-availability single-reply-date diagnostic with 328,695 events.

We add group indicators and four ESI-by-group terms to equation~\eqref{eq:rpr}, first with baseline controls and then with both tones. Table~\ref{tab:lag_regressions} gives total ESI slopes. All are positive, but joint equality tests have $p=0.1213$ and $0.1247$. Within-group portfolios in Table~\ref{tab:lag_portfolios} likewise do not establish a monotonic benefit from faster replies. All 40 paired alpha comparisons across groups, weights and factors have unadjusted $p>0.05$. This does not prove speed irrelevant: estimates may be imprecise and speed is endogenous to questions and firms. Individual group stars cannot substitute for tests of slope differences.
\begin{table}[H]
\centering\footnotesize
\caption{Reply Lag Distribution and Historical O2C Slopes}\label{tab:lag_regressions}
\begin{threeparttable}
\renewcommand{\arraystretch}{1.12}
\begin{tabular*}{\textwidth}{@{\extracolsep{\fill}}lrrrr@{}}
\toprule
Maximum lag & Firm days & Share (\%) & Baseline & Also both tones \\
\midrule
Same day & 100,022 & 11.56 & 3.05 (3.75) & 3.44 (3.60) \\
1 day & 190,770 & 22.05 & 2.80 (4.47) & 3.18 (4.08) \\
2 days & 92,125 & 10.65 & 1.92 (2.15) & 2.30 (2.24) \\
3--5 days & 187,013 & 21.61 & 0.88 (1.32) & 1.26 (1.54) \\
More than 5 days & 295,379 & 34.14 & 1.43 (2.88) & 1.83 (2.69) \\
\bottomrule
\end{tabular*}
\begin{tablenoteblock}
The maximum calendar-day interval among recorded replies in each original company-question-date cohort defines the group. All matched replies must have valid nonnegative lags; no single-reply-date restriction is imposed. The common grouped sample is 865,309 firm days, not the 328,695 observations of the separate public-availability diagnostic. Slopes are total ESI slopes in bps per frozen full-regression-pool standard deviation, with $t$ statistics in parentheses. The pooled regression includes group indicators, ESI-by-group terms, eleven controls and firm/date effects; errors are clustered by firm/date. Joint slope-equality $p$ values are 0.1213 and 0.1247 (Holm-adjusted 0.2425 for both).
\end{tablenoteblock}
\end{threeparttable}
\end{table}

\begin{table}[H]
\centering\footnotesize
\caption{ESI Portfolios within Reply Lag Groups}\label{tab:lag_portfolios}
\begin{threeparttable}
\renewcommand{\arraystretch}{1.12}
\begin{tabular*}{\textwidth}{@{\extracolsep{\fill}}lrrrr@{}}
\toprule
Maximum lag & EW q5 $\alpha$ & Dates & VW q5 $\alpha$ & Dates \\
\midrule
Same day & 4.47 (2.86) & 3,570 & 5.28 (2.22) & 3,570 \\
1 day & 2.91 (1.67) & 3,537 & 3.12 (1.38) & 3,537 \\
2 days & 4.12 (1.41) & 3,250 & 2.70 (0.82) & 3,250 \\
3--5 days & 3.31 (2.10) & 3,529 & 3.07 (1.48) & 3,529 \\
More than 5 days & 2.50 (1.86) & 3,634 & 5.29 (2.79) & 3,634 \\
\bottomrule
\end{tabular*}
\begin{tablenoteblock}
The original O2C sorting and frozen portfolio ESI rules apply separately within each maximum-calendar-lag group. Alphas are in bps and HAC(5) $t$ statistics are in parentheses. Dates are own-strategy dates, so a comparison of separate estimates is not a test of their difference. All 40 paired alpha differences across five groups, two weights and two factor benchmarks have unadjusted $p>0.05$ (minimum 0.2015); the complete results are retained in the repository. The U.S. q5 source and historical timing limitations remain applicable.
\end{tablenoteblock}
\end{threeparttable}
\end{table}

